% Part: proof-theory
% Chapter: cut-elimination
% Section: ce-largest

\documentclass[../../../include/open-logic-section]{subfiles}

\begin{document}

\olfileid{pt}{cut}{inv}

\olsection{Ngilangi Cut kang Pangkaté Paling Gedhé}

Bukti \olref[seq][inv]{lem:G3c-invert} nuduhaké yèn yèn
\Log{G3c} !!{prove}s dudutan aturan logis
(saliyané \RightR{\lexists} lan \LeftR{\lforall}),
saben premis aturan mau uga bisa !!{prove}s tanpa nambah
!!{height} buktiné. Kita bisa nuduhaké pratelan kang luwih
kuwat: iki uga bener kanggo $\Log{G3c} + \CutCS$.

Kaya sadurungé, \emph{pangkat cut} saka inferènsi
\CutCS{} kita tetepaké minangka ambane !!{formula}
kang di-cut.

\begin{lem}\ollabel{lem:inv-G3c-cut}
Yèn $\Log{G3c} + \CutCS$ !!{prove}s dudutan aturan logis
~$R$ saliyané \RightR{\lexists} utawa \LeftR{\lforall}
nganggo !!a{proof}~$\pi$, saben premisé uga bisa
!!{prove}s nganggo !!{proof}~$\pi'$
(utawa !!{proof}s~$\pi'$, $\pi''$ yèn aturané
nduwèni rong premis). Kajaba iku,
(a)~$\pheight{\pi'}\le\pheight{\pi}$
(lan $\pheight{\pi''}\le\pheight{\pi}$), sarta
(b)~cacah inferènsi \CutCS{} lan pangkaté ing
~$\pi'$ (lan $\pi''$) padha karo kang ana ing~$\pi$.
\end{lem}

\begin{proof}
Delengen bukti
\cref{pt:seq:inv:lem:G3c-invert,pt:seq:inv:lem:invert-quant}.
Bukti mau nganggo induksi tumrap~$\pheight{\pi}$.
Ing kahanan dhasar, sawijining aksioma diganti aksioma liya;
mula syarat (a) lan~(b) langsung kacukupan.
Yèn aturan pungkasané~$\pi$ yaiku~$R$,
sub-!!{proof}s langsungé yaiku !!{proof}s~$\pi_i$
kang dikarepaké, lan (a) sarta~(b) uga kacukupan.
Ing kahanan liya, hipotesis induksi ditrapaké marang
sub-!!{proof}s langsung saka~$\pi$, yaiku kang nuju
premis-premis inferènsi pungkasan~$R$; banjur aturan
$R$ kang padha ditrapaké marang asilé.
Syarat (a) lan~(b) bener kanggo sub-!!{proof}s langsung
saka bukti anyar, mula uga bener kanggo sakabèhé bukti.
Isih kudu dipriksa kahanan nalika inferènsi pungkasané
yaiku~\CutCS. Kaya mau, cukup siji tuladha: anggepen
$\pi$ iku !!a{proof} kanggo dudutan
$!B \land !C, \Gamma \Sequent \Delta$ saka
\LeftR{\land}, lan pungkasané nganggo~\CutCS:
\[
\AxiomC{}
\RightLabel{$\pi_1$}
\Deduce$!B \land !C, \Gamma \fCenter \Delta, !A$
\AxiomC{}
\RightLabel{$\pi_2$}
\Deduce$!A, !B \land !C, \Gamma \fCenter \Delta$
\RightLabel{\CutCS}
\BinaryInf$!B \land !C, \Gamma \fCenter \Delta$
\DisplayProof
\]
Hipotesis induksi ditrapaké marang $\pi_1$ lan~$\pi_2$
(kang uga !!{proof}s kanggo tuladha dudutan
$\LeftR{\land}$), lan ngasilaké !!{proof}s
$\pi_1'$ lan~$\pi_2'$ kanggo
$!B, !C, \Gamma \fCenter \Delta, !A$ lan
$!A, !B, !C, \Gamma \fCenter \Delta$ miturut
urutané. Kaloroné digandhèng nganggo \CutCS{}
kanggo ngolèhké~$\pi'$:
\[
\AxiomC{}
\RightLabel{$\pi_1'$}
\Deduce$!B, !C, \Gamma \fCenter \Delta, !A$
\AxiomC{}
\RightLabel{$\pi_2'$}
\Deduce$!A, !B, !C, \Gamma \fCenter \Delta$
\RightLabel{\CutCS}
\BinaryInf$!B, !C, \Gamma \fCenter \Delta$
\DisplayProof
\]
Cetha yèn (a) lan~(b) kacukupan.
\end{proof}

Élinga yèn iki ora bakal lumaku yèn kita nganggo
\Cut{} tinimbang \CutCS, amarga dudutan !!{proof}
pungkasan bakal nduwèni rong salinan
$!B, !C, \Gamma$ ing antesèden lan rong salinan
$\Delta$ ing suksèden. Kita kudu migunakaké
admisibilitas kontraksi, nanging bukti
\olref[seq][inv]{prop:G3c-cont-adm} migunakaké lema
inversibilitas. Iki bakal dadi argumèn mbunder.

Lema \cref{pt:cut:inv:lem:inv-G3c-cut} bisa digunakaké
kanggo bukti alternatif eliminasi cut. Ing bukti iki,
kita ora ngilangi inferènsi \CutCS{} paling dhuwur
siji-siji, nanging ngilangi inferènsi \CutCS{}
kang pangkaté maksimal. Wiwitané ana !!a{proof}~$\pi$
ing $\Log{G3c} + \CutCS$; sawijining cut ing endi waé
dibusak (bisa uga diganti cut kang pangkaté luwih cilik),
lan iki diulang nganti ora ana cut. Syaraté mung siji:
ing sadhuwuré cut kang digarap, ora ana cut kang
pangkaté padha utawa luwih gedhé.

\begin{lem}\ollabel{lem:max-cut-red-G3c} Yèn $\pi$
iku !!a{proof} ing $\Log{G3c}+\CutCS$ kang pungkasané
nganggo inferènsi \CutCS{} pangkat~$n$, lan saliyané
mung nganggo aturan~$\Log{G3c}$ sarta inferènsi
\CutCS{} pangkat $<n$, mula ana bukti~$\pi'$
kanggo sekèn pungkasan kang padha ing
$\Log{G3c} + \CutCS$, kanthi saben inferènsi
\CutCS{} pangkaté~$<n$.
\end{lem}

\begin{proof}
!!{proof}~$\pi$ pungkasané kaya mangkéné:
\[
\AxiomC{}
\RightLabel{$\pi_1$}
\Deduce$\Gamma \fCenter \Delta, !A$
\AxiomC{}
\RightLabel{$\pi_2$}
\Deduce$!A, \Gamma \fCenter \Delta$
\RightLabel{\CutCS}
\BinaryInf$\Gamma \fCenter \Delta$
\DisplayProof
\]
Kita bedakaké kahanané manut wujudé~$!A$.

Anggepen $!A$ atomik. Miturut syarat yèn kabèh inferènsi
\CutCS{} ing~$\pi_1$ lan $\pi_2$ kudu pangkaté $<
\depth{!A} = 0$, ora ana inferènsi \CutCS{} ing~$\pi_1$
utawa~$\pi_2$. Kita nuduhaké yèn
$\Gamma \Sequent \Delta$ nduwèni bukti tanpa cut
kanthi induksi tumrap $\pheight{\pi_2}$.
Yèn $\pheight{\pi_2} = 0$, banjur $!A,
\Gamma \Sequent \Delta$ iku aksioma. Yèn $!A$ dudu
formula utama, ana $!C$ atomik kang dadi !!a{element} ing~$\Gamma$
lan $\Delta$, mula $\Gamma \Sequent \Delta$ dhéwé
iku aksioma. Yèn $!A$ formula utama, banjur
$\Delta = \Delta', !A$. Ing kahanan iki, $\pi_1$ pungkasané
$\Gamma \Sequent \Delta', !A, !A$. Kita ngolèhké
bukti tanpa cut kanggo $\Gamma \Sequent \Delta', !A$
kanthi ngetrapaké \olref[seq][inv]{prop:G3c-cont-adm}
(admisibilitas kontraksi). Yèn $\pheight{\pi_2} > 0$,
bukti iki pungkasané nganggo aturan~$R$. Jupuken salah siji
premis aturan mau: wujudé $!A, \Gamma',
\Pi \Sequent \Delta', \Lambda$, ing ngendi $\Pi$ lan $\Lambda$
iku !!{formula}s aktif saka~$R$. Kanthi ngetrapaké
\olref[seq][adm]{prop:weak-G3c-adm}
marang~$\pi_1$ lan $\pi_2$, kita ngolèhké !!{proof}s
kanggo $\Gamma, \Gamma',
\Pi \Sequent \Delta, \Delta', \Lambda, !A$ lan $!A, \Gamma, \Gamma',
\Pi \Sequent \Delta, \Delta', \Lambda$; hipotesis induksi
ditrapaké marang kaloroné. Mula kita ngolèhké
!!a{proof} kanggo $\Gamma, \Gamma', \Pi
\Sequent \Delta, \Delta', \Lambda$ kanggo saben premis~$R$.
Ngetrapaké $R$ ngasilaké $\Gamma, \Gamma \Sequent \Delta, \Delta$,
lan \olref[seq][inv]{prop:G3c-cont-adm} maringi bukti tanpa cut
kanggo $\Gamma \Sequent \Delta$.

Yèn $!A$ ora atomik, kita bedakaké kahanané manut !!{main
operator} saka~$!A$. Ing saben kahanan, kita ngetrapaké
\olref{lem:inv-G3c-cut} kanggo ngolèhké !!{proof}s
premis-premis aturan kiwa lan tengen kanthi $!A$ minangka
formula utama. Banjur kita nerusaké kaya ing kahanan~E
saka bukti \olref[top]{lem:cut-adm-G3c}.
\begin{enumerate}
  \item $!A \ident \lnot !B$. !!{proof}s $\pi_1$ lan $\pi_2$
  pungkasané ana ing $\Gamma \Sequent \Delta, \lnot !B$ lan
  $\lnot !B, \Gamma \Sequent \Delta$. Miturut
  \olref{lem:inv-G3c-cut}, ana !!{proof}s $\pi_1'$
  lan $\pi_2'$ kanggo $!B, \Gamma \Sequent
  \Delta$ lan $\Gamma \Sequent \Delta, !B$.
  !!{proof}~$\pi'$ iki:
  \[
  \AxiomC{}
  \RightLabel{$\pi_2'$}
  \Deduce$\Gamma \fCenter \Delta, !B$
  \AxiomC{}
  \RightLabel{$\pi_1'$}
  \Deduce$!B, \Gamma \fCenter \Delta$
  \RightLabel{\CutCS}
  \BinaryInf$\Gamma \fCenter \Delta$
  \DisplayProof
  \]
  \item $!A \ident (!B \lor !C)$. !!{proof}s $\pi_1$ lan $\pi_2$
  pungkasané ana ing $\Gamma \Sequent \Delta, !B \lor !C$ lan
  $!B \lor !C, \Gamma \Sequent \Delta$. Miturut
  \olref{lem:inv-G3c-cut} (inversi kanggo
  \RightR{\lor} kang ditrapaké marang~$\pi_1$), ana
  !!a{proof} $\pi_1'$ kanggo
  $\Gamma \Sequent \Delta, !B, !C$. Miturut
  \olref{lem:inv-G3c-cut} (inversi kanggo \LeftR{\lor}
  kang ditrapaké marang~$\pi_2$), ana !!a{proof}
  $\pi_2'$ kanggo $!B, \Gamma \Sequent \Delta$ lan
  !!a{proof} $\pi_2''$ kanggo
  $!C, \Gamma \Sequent \Delta$. Kanthi ngetrapaké
  \olref[seq][adm]{prop:weak-G3c-adm} marang $\pi_2''$,
  kita uga ngolèhké !!a{proof}~$\pi_2'''$ kanggo
  $!C, \Gamma \Sequent \Delta, !B$. !!{proof}~$\pi'$ iki:
  \[
  \AxiomC{}
  \RightLabel{$\pi_1'$}
  \Deduce$\Gamma \fCenter \Delta, !B, !C$
  \AxiomC{}
  \RightLabel{$\pi_2'''$}
  \Deduce$!C, \Gamma \fCenter \Delta, !B$
  \RightLabel{\CutCS}
  \BinaryInf$\Gamma \fCenter \Delta, !B$
  \AxiomC{}
  \RightLabel{$\pi_2'$}
  \Deduce$!B, \Gamma \fCenter \Delta$
  \RightLabel{\CutCS}
  \BinaryInf$\Gamma \fCenter \Delta$
  \DisplayProof
  \]
  \item $!A \ident (!B \land !C)$. Latihan.
  \item $!A \ident (!B \lif !C)$. !!{proof}s $\pi_1$ lan $\pi_2$
  pungkasané ana ing $\Gamma \Sequent \Delta, !B \lif !C$ lan
  $!B \lif !C, \Gamma \Sequent \Delta$. Miturut
  \olref{lem:inv-G3c-cut} (inversi kanggo
  \RightR{\lif} kang ditrapaké marang~$\pi_1$), ana
  !!a{proof} $\pi_1'$ kanggo
  $!B, \Gamma \Sequent \Delta, !C$. Miturut
  \olref{lem:inv-G3c-cut} (inversi kanggo
  \LeftR{\lif} kang ditrapaké marang~$\pi_2$), ana
  !!a{proof} $\pi_2'$ kanggo $\Gamma \Sequent \Delta, !B$
  lan !!a{proof} $\pi_2''$ kanggo
  $!C, \Gamma \Sequent \Delta$. Kanthi ngetrapaké
  \olref[seq][adm]{prop:weak-G3c-adm} marang $\pi_2''$,
  kita uga ngolèhké !!a{proof}~$\pi_2'''$ kanggo
  $!C, !B, \Gamma \Sequent \Delta$. !!{proof}~$\pi'$ iki:
  \[
  \AxiomC{}
  \RightLabel{$\pi_2'$}
  \Deduce$\Gamma \fCenter \Delta, !B$
  \AxiomC{}
  \RightLabel{$\pi_1'$}
  \Deduce$!B, \Gamma \fCenter \Delta, !C$
  \AxiomC{}
  \RightLabel{$\pi_2'''$}
  \Deduce$!C, !B, \Gamma \fCenter \Delta$
  \RightLabel{\CutCS}
  \BinaryInf$!B, \Gamma \fCenter \Delta$
  \RightLabel{\CutCS}
  \BinaryInf$\Gamma \fCenter \Delta$
  \DisplayProof
  \]
  \item $!A \ident \lforall[x][!B(x)]$. !!{proof}s $\pi_1$ lan
  $\pi_2$ pungkasané ana ing
  $\Gamma \Sequent \Delta, \lforall[x][!B(x)]$ lan
  $\lforall[x][!B(x)], \Gamma \Sequent \Delta$. Miturut
  \olref{lem:inv-G3c-cut}, ana !!a{proof}~$\pi_1'$ kanggo
  $\Gamma \Sequent \Delta, !B(c)$.

  Saiki gatekna !!{proof}~$\pi_2$. Bukti iki pungkasané ana ing
  $\lforall[x][!B(x)], \Gamma \Sequent \Delta$. Saka sekèn
  pungkasan ing~$\pi_2$, lumakua munggah lan tandhanana
  saben kedadéan $\lforall[x][!B(x)]$ ing sisih kiwa sekèn kang
  ``nuju'' menyang kedadéan $\lforall[x][!B(x)]$ ing sekèn
  pungkasan~$\pi_2$, yaiku:
  (a)~Tandhanana $\lforall[x][!B(x)]$ ing sekèn pungkasan~$\pi_2$.
  (b)~Yèn $\lforall[x][!B(x)]$ ana ing kontèks konklusi aturan
  lan wis ditandhani, tandhanana kedadéan kang cocog ing
  premis utawa premis-premis aturan mau.
  (c)~Yèn $\lforall[x][!B(x)]$ iku formula utama ing konklusi
  aturan \LeftR{\lforall} lan wis ditandhani, tandhanana
  kedadéan aktif kang cocog saka $\lforall[x][!B(x)]$ lan $!B(t)$
  ing premisé.

  Saiki gatekna kabèh kedadéan $\lforall[x][!B(x)]$ kang
  ditandhani. Ora ana kang aktif ing aturan saliyané
  \LeftR{\lforall} (mung !!{formula}s aktif saka
  \LeftR{\lforall} kang kena tandha). Busaken kabèh
  kedadéan $\lforall[x][!B(x)]$ kang ditandhani saka $\pi_2$.
  Yèn asilé tetep !!{proof} kang bener, kedadéan
  $\lforall[x][!B(x)]$ kang ditandhani mau ora tau aktif;
  kabèh asalé langsung saka aksioma. !!{proof} iki pungkasané
  ana ing~$\Gamma \Sequent \Delta$ lan bisa ngganti $\pi$.
  Mangkono siji cut pangkat maksimal wis kabusak.

  Yèn ora, asilé dudu bukti kang bener amarga kita wis
  mbusak !!{formula}s $\lforall[x][!B(x)]$ kang aktif ing
  sawatara inferènsi \LeftR{\lforall}. Inferènsi-inferènsi
  mau dadi ``inferènsi'' kanthi wujud
  \[
  \AxiomC{}
  \RightLabel{$\pi_t$}
  \Deduce$!B(t), \Pi \fCenter \Lambda$
  \RightLabel{\LeftR{\lforall}}
  \UnaryInf$\Pi \fCenter \Lambda$
  \DisplayProof
  \]
  Saben ``inferènsi'' mangkono kang paling dhuwur diganti karo
  \[
  \AxiomC{}
  \RightLabel{$\Subst{\pi_1'}{t}{c}[\Pi \Sequent \Lambda]$}
  \Deduce$\Gamma, \Pi \fCenter \Delta, \Lambda, !B(t)$
  \AxiomC{}
  \RightLabel{$\pi_t[\Gamma \Sequent \Delta]$}
  \Deduce$!B(t), \Gamma, \Pi \fCenter \Delta, \Lambda$
  \RightLabel{\CutCS}
  \BinaryInf$\Gamma, \Pi \fCenter \Delta, \Lambda$
  \DisplayProof
  \]
  banjur tambahna $\Gamma$ ing sisih kiwa lan $\Delta$ ing sisih
  tengen saben sekèn ing sangisoré. Balènana nganti kabèh
  ``inferènsi'' \LeftR{\lforall} kang premisé ngandhut
  $!B(t)$ bertandha diganti inferènsi \CutCS{} tumrap
  sawijining~$!B(t)$. Sekèn pungkasan saka !!{proof} asilé
  (kang saiki dadi !!{proof} kang bener) wujudé
  $\Gamma, \dots, \Gamma \Sequent \Delta, \dots, \Delta$.
  Trapna admisibilitas kontraksi supaya dadi !!a{proof} kanggo
  $\Gamma \Sequent \Delta$, banjur gantenana $\pi$ nganggo bukti
  iku. Siji cut tumrap $\lforall[x][!B(x)]$ wis diganti
  (mbokmenawa akèh) cut tumrap $!B(t)$, nanging pangkaté
  kabèh luwih cilik. Cut kang wis ana ing $\pi_1$ utawa
  $\pi_2$ tetep ana, nanging pangkaté uga kabèh luwih cilik.
\end{enumerate}
\end{proof}

\begin{prob}
Priksanen kahanan $!A \ident (!B \land !C)$ ing bukti
\olref{lem:max-cut-red-G3c}. Tegesé, tuduhna yèn ana
!!a{proof}~$\pi$ kang pungkasané kaya mangkéné:
\[
\AxiomC{}
\RightLabel{$\pi_1$}
\Deduce$\Gamma \fCenter \Delta, !B \land !C$
\AxiomC{}
\RightLabel{$\pi_2$}
\Deduce$!B \land !C, \Gamma \fCenter \Delta$
\RightLabel{\CutCS}
\BinaryInf$\Gamma \fCenter \Delta$
\DisplayProof
\]
lan $\pi_1$ sarta $\pi_2$ mung ngemot cut kang pangkaté
$<\depth{!B \land !C}$, banjur ana !!a{proof}~$\pi'$ kanggo
$\Gamma \Sequent \Delta$ kang mung ngemot cut kang pangkaté
$<\depth{!B \land !C}$.
\end{prob}

% \begin{prob}
% Verify the case where $!A \ident \lexists[x][!B(x)]$ in the proof of
% \olref{lem:inv-G3c-cut}.
% \end{prob}


\olref{lem:max-cut-red-G3c} nuduhaké yèn inferènsi \CutCS{}
kang pangkaté paling gedhé ing !!a{proof} bisa diganti
inferènsi-inferènsi \CutCS{} kang pangkaté luwih cilik.
Lema iki bisa digunakaké manèh kanggo nuduhaké yèn
pirang-pirang inferènsi \CutCS{} ing !!a{proof} bisa
diilangi kanthi ngetrapaké lema mau kanthi urut marang
sub-!!{proof}s paling dhuwur kang pungkasané ana ing
\CutCS{} pangkat maksimal.

\begin{thm}\ollabel{thm:cut-elim-G3c-largest}
Saben !!{proof} $\pi$ ing $\Log{G3c} + \CutCS$
bisa diowahi dadi bukti ing~\Log{G3c}.
\end{thm}

\begin{proof}
  Sepisan, kita buktèkaké yèn kabèh cut kang pangkaté maksimal
  ing~$\pi$ bisa diilangi. Upama $d$ iku pangkat maksimal cut
  kang ana ing~$\pi$, lan $n$ cacah cut pangkat~$d$.
  Buktiné nganggo induksi tumrap~$n$. Yèn $n=0$, ora ana
  kang perlu dibuktèkaké. Yèn $n>0$, pilihen cut pangkat~$d$
  kang paling dhuwur, lan upama $\pi_1$ sub-!!{proof} kang
  pungkasané ana ing cut iku. Miturut
  \olref{lem:max-cut-red-G3c}, ana !!a{proof}~$\pi_1'$
  kanggo sekèn pungkasan kang padha kanthi kabèh cut
  pangkat $<d$. Gantenana $\pi_1$ nganggo $\pi_1'$ ing~$\pi$.
  Asilé !!a{proof} kang ngemot $n-1$ cut pangkat~$d$,
  mula hipotesis induksi bisa ditrapaké.

  Saiki asilé ndhèrèk saka induksi tumrap~$d$. Yèn $d=0$,
  kabèh cut atomik, lan miturut asil sadurungé kabèh bisa
  diilangi. Yèn $d>0$, asil sadurungé maringi bukti kang
  kabèh cut pangkat $d$ wis diganti cut pangkat~$<d$.
  Amarga $d$ iku pangkat maksimal cut ing~$\pi$, saiki
  kita nduwèni bukti kang pangkat maksimal cuté $<d$,
  mula hipotesis induksi bisa ditrapaké.
\end{proof}

\end{document}
